\section{Conclusion}
\label{sec:conclusion}

We investigated whether \futurelens-derived prediction metrics can detect \noncomp phrases memorized by LLMs.
Using \gptxl and two compositionality-annotated datasets, we find that \noncomp phrases are significantly more predictable from intermediate hidden states than \comp ones.
On the \farahmand noun compound dataset, the best single metric (next-token probability) achieves ROC AUC = 0.643, and prediction metrics correlate with human compositionality scores at Spearman $r = 0.23$ ($p < 0.0001$).
On the \idiomem idiom dataset, the effect is dramatically larger (Cohen's $d = 1.85$), confirming that memorized \noncomp units leave a clear signature in model internals.

We find that within-phrase prediction signals are more robust than from-first-token prediction, and that the \noncomp advantage emerges at layer 6 and grows through the final layers of the model.
These results support the hypothesis that \futurelens can lower-bound the set of \noncomp phrases an LLM has memorized, though the signal is moderate in isolation.

Future work should train dedicated \futurelens probes (rather than relying on the \logitlens), control for phrase frequency, combine forward prediction with backward erasure signals~\citep{feucht2024token}, and replicate on larger models.
The prediction metrics we evaluate here are best suited as components in a multi-signal non-compositionality detection system.
